\lecture{23}{Метод множителей}

\sourcevideo
    {https://www.youtube.com/playlist?list=PLOzRYVm0a65fhKDS8cYIC7YCuVGJ4FOJx}
    {Week 7: Lecture 23: Method of Multipliers}

Рассмотрим задачу
\[
    \min_{x\in\mathbb R^n}f(x)
    \qquad
    \text{при условии}
    \qquad
    h(x)=0,
\]
где
\[
    h\colon\mathbb R^n\to\mathbb R^m.
\]
Дополненная функция Лагранжа имеет вид
\[
    \mathcal L_c(x,\nu)
    =
    f(x)
    +
    \nu^{\mathsf T}h(x)
    +
    \frac c2\lVert h(x)\rVert^2,
    \qquad
    c>0.
\]
Метод множителей чередует минимизацию по прямой переменной и
коррекцию двойственного множителя по невязке ограничения.

\section{Схема метода и стационарность}

При заданных \(\nu^k\) и \(c_k>0\) выполняется прямой шаг
\[
    x^{k+1}
    \in
    \operatorname*{argmin}_{x\in\mathbb R^n}
    \mathcal L_{c_k}(x,\nu^k),
\]
после чего множитель обновляется по правилу
\[
    \nu^{k+1}
    =
    \nu^k+c_kh(x^{k+1}).
\]
В лекции дополнительно рассматривается увеличение штрафа
\[
    c_{k+1}=\beta c_k,
    \qquad
    \beta>1,
\]
и упоминаются значения \(\beta\in[5,10]\). Такое быстрое увеличение
может ухудшать обусловленность, поэтому на практике часто используют
умеренный рост или фиксируют \(c_k=c\).

Прямая подзадача является безусловной:
\[
    \min_x
    \left\{
        f(x)
        +
        (\nu^k)^{\mathsf T}h(x)
        +
        \frac{c_k}{2}\lVert h(x)\rVert^2
    \right\}.
\]
Её можно решать градиентным, инерционным, ньютоновским или
квазиньютоновским методом. В теоретической записи минимизатор
находится точно, а в реализации допускается приближённое решение.

Предположим сначала, что прямой шаг точен. Тогда
\[
    \nabla_x\mathcal L_{c_k}(x^{k+1},\nu^k)=0.
\]
Поскольку
\[
\begin{aligned}
    \nabla_x\mathcal L_{c_k}(x,\nu^k)
    ={}&
    \nabla f(x)
    +
    J_h(x)^{\mathsf T}\nu^k
    \\
    &+
    c_kJ_h(x)^{\mathsf T}h(x),
\end{aligned}
\]
получаем
\[
    \nabla f(x^{k+1})
    +
    J_h(x^{k+1})^{\mathsf T}\nu^{k+1}
    =0.
\]
Таким образом, прямой шаг и обновление множителя автоматически дают
стационарность обычной функции Лагранжа. Для выполнения KKT-условий
остаётся обеспечить допустимость
\[
    h(x^{k+1})\to0.
\]
Из двойственного обновления непосредственно следует
\[
    h(x^{k+1})
    =
    \frac{\nu^{k+1}-\nu^k}{c_k}.
\]
При фиксированном \(c_k=c>0\) достаточно затухания приращений
множителя. При ограниченной последовательности \(\{\nu^k\}\) и
\(c_k\to\infty\) невязка также стремится к нулю при дополнительных
условиях существования предельных точек.

\section{Предельные точки и неточный прямой шаг}

При стандартных предпосылках сходимости, если \(f\) и \(h\)
непрерывны, допустимое множество непусто, прямые подзадачи решаются
глобально, последовательность множителей ограничена, а
\(c_k\to\infty\), то при стандартных предпосылках компактности всякая
предельная точка \(\{x^k\}\) является глобальным решением исходной
задачи. Возрастающий штраф не позволяет предельной точке сохранять
ненулевую невязку.

Для невыпуклой задачи глобальное решение прямой подзадачи обычно
недоступно. Если внутренний алгоритм находит лишь локальный минимум,
то в общем случае сохраняются только локальные гарантии.

При приближённой минимизации вместо точного равенства используется
условие
\[
    \left\lVert
        \nabla_x
        \mathcal L_{c_k}(x^{k+1},\nu^k)
    \right\rVert
    \le
    \varepsilon_k,
    \qquad
    \varepsilon_k\to0.
\]
На ранних внешних итерациях допустим грубый внутренний допуск, а по
мере сходимости точность увеличивается. Для сходимости требуется согласовать скорость убывания \(\varepsilon_k\) с выбором \(c_k\). При полном ранге якобиана
ограничения и подходящем убывании ошибки можно получить также
сходимость приближённых множителей
\[
    \nu^k+c_kh(x^{k+1})\to\nu^\star.
\]

Большой штраф улучшает допустимость, но может сделать внутреннюю
задачу плохо обусловленной. Для линейного ограничения
\[
    h(x)=Ax-b
\]
Гессиан прямой подзадачи содержит
\[
    \nabla^2f(x)+c_kA^{\mathsf T}A.
\]
Кривизна растёт только в направлениях, на которые действует \(A\),
и не меняется вдоль \(\ker A\). Поэтому увеличение \(c_k\) может
резко увеличить число обусловленности. Практически используют тёплый
запуск из предыдущей прямой точки, предобусловливание, ньютоновские
или квазиньютоновские методы и не увеличивают штраф без необходимости.

\section{Выпуклый квадратичный пример}

Рассмотрим задачу
\[
    \min_{x_1,x_2}
    \frac12(x_1^2+x_2^2)
    \qquad
    \text{при условии}
    \qquad
    x_1=1.
\]
Её решение и оптимальный множитель равны
\[
    x^\star=
    \begin{pmatrix}
        1\\
        0
    \end{pmatrix},
    \qquad
    \nu^\star=-1.
\]
Дополненная функция Лагранжа имеет вид
\[
\begin{aligned}
    \mathcal L_c(x,\nu)
    ={}&
    \frac12(x_1^2+x_2^2)
    +
    \nu(x_1-1)
    \\
    &+
    \frac c2(x_1-1)^2.
\end{aligned}
\]
Условия минимума по прямым переменным дают
\[
    x_1^{k+1}
    =
    \frac{c_k-\nu^k}{1+c_k},
    \qquad
    x_2^{k+1}=0.
\]
Следовательно,
\[
    x_1^{k+1}-1
    =
    -\frac{\nu^k+1}{1+c_k}.
\]
После двойственного обновления
\[
    \nu^{k+1}
    =
    \nu^k+c_k(x_1^{k+1}-1)
\]
получаем точную рекурсию
\[
    \nu^{k+1}-\nu^\star
    =
    \frac{\nu^k-\nu^\star}{1+c_k}.
\]
Для любого \(c_k>0\) коэффициент лежит в \((0,1)\), поэтому ошибка
множителя монотонно сокращается. При постоянном \(c_k=c\)
\[
    |\nu^k-\nu^\star|
    =
    \frac{|\nu^0-\nu^\star|}{(1+c)^k}.
\]
Рост штрафа до бесконечности в этом примере не нужен.

Гессиан прямой подзадачи равен
\[
    \nabla^2_{xx}\mathcal L_c
    =
    \begin{pmatrix}
        1+c&0\\
        0&1
    \end{pmatrix},
\]
поэтому
\[
    \kappa=1+c.
\]
Большое \(c\) уменьшает внешний коэффициент сжатия \(1/(1+c)\), но
одновременно ухудшает внутреннюю обусловленность. Это основной
компромисс метода множителей.

\section{Невыпуклый квадратичный пример}

Теперь рассмотрим задачу
\[
    \min_{x_1,x_2}
    \left\{
        -\frac12x_1^2+\frac12x_2^2
    \right\}
    \qquad
    \text{при условии}
    \qquad
    x_1=1.
\]
Хотя целевая функция невыпукла по \(x_1\), ограничение фиксирует эту
координату. Поэтому
\[
    x^\star=
    \begin{pmatrix}
        1\\
        0
    \end{pmatrix},
    \qquad
    \nu^\star=1.
\]
Дополненная функция Лагранжа равна
\[
\begin{aligned}
    \mathcal L_c(x,\nu)
    ={}&
    -\frac12x_1^2
    +
    \frac12x_2^2
    +
    \nu(x_1-1)
    \\
    &+
    \frac c2(x_1-1)^2.
\end{aligned}
\]
Её Гессиан имеет вид
\[
    \nabla^2_{xx}\mathcal L_c
    =
    \begin{pmatrix}
        c-1&0\\
        0&1
    \end{pmatrix}.
\]
Следовательно, прямая подзадача сильно выпукла только при
\[
    c>1.
\]
Для \(c_k>1\) прямой шаг равен
\[
    x_1^{k+1}
    =
    \frac{c_k-\nu^k}{c_k-1},
    \qquad
    x_2^{k+1}=0,
\]
а невязка имеет вид
\[
    x_1^{k+1}-1
    =
    -\frac{\nu^k-1}{c_k-1}.
\]
После двойственного обновления получаем
\[
    \nu^{k+1}-\nu^\star
    =
    -\frac{\nu^k-\nu^\star}{c_k-1}.
\]
Знак минус означает
чередование направления ошибки. Для сжатия по модулю необходимо
\[
    \frac1{c_k-1}<1,
    \qquad
    \text{то есть}
    \qquad
    c_k>2.
\]
При \(1<c_k<2\) прямая подзадача уже сильно выпукла, но внешняя
итерация расходится; при \(c_k=2\) ошибка только меняет знак; при
\(c_k>2\) возникает сжатие. Улучшение кривизны внутренней задачи само
по себе поэтому не гарантирует сходимость всего метода.

В общих задачах порог заранее может быть неизвестен. Штраф начинают с
умеренного значения и увеличивают только при недостаточном уменьшении
невязки. Слишком малый \(c_k\) слабо обеспечивает допустимость и может
не устранить отрицательную кривизну, а слишком большой ухудшает
обусловленность, уменьшает допустимый внутренний шаг и усиливает
численные ошибки.

\section{Неравенства и проекционное обновление}

Для задачи
\[
    \min_x f(x)
\]
при условиях
\[
    h(x)=0,
    \qquad
    g_j(x)\le0,
    \quad
    j=1,\ldots,r,
\]
каждое неравенство можно заменить равенством
\[
    g_j(x)+z_j^2=0,
    \qquad
    z_j\in\mathbb R.
\]
Такое равенство разрешимо по \(z_j\) тогда и только тогда, когда
\(g_j(x)\le0\). В векторной форме используется
\[
    g(x)+z\circ z=0.
\]
Дополненная функция Лагранжа принимает вид
\[
\begin{aligned}
    \mathcal L_c(x,z,\nu,\lambda)
    ={}&
    f(x)
    +
    \nu^{\mathsf T}h(x)
    +
    \frac c2\lVert h(x)\rVert^2
    \\
    &+
    \lambda^{\mathsf T}(g(x)+z\circ z)
    \\
    &+
    \frac c2
    \lVert g(x)+z\circ z\rVert^2.
\end{aligned}
\]
Прямой шаг выполняется по \((x,z)\), а множители обновляются как
\[
    \nu^{k+1}
    =
    \nu^k+c_kh(x^{k+1}),
\]
\[
\begin{aligned}
    \lambda^{k+1}
    ={}&
    \lambda^k
    +
    c_k
    \bigl(
        g(x^{k+1})
        +
        z^{k+1}\circ z^{k+1}
    \bigr).
\end{aligned}
\]

Переменные \(z_j\) можно исключить. Для одной компоненты минимизируется
по \(t=z_j^2\ge0\) выражение
\[
    \lambda_j(g_j(x)+t)
    +
    \frac c2(g_j(x)+t)^2.
\]
После исключения получается проекционное обновление
\[
    \lambda_j^{k+1}
    =
    \max\left\{
        0,
        \lambda_j^k+c_kg_j(x^{k+1})
    \right\},
\]
или в векторной форме
\[
    \lambda^{k+1}
    =
    \bigl[
        \lambda^k+c_kg(x^{k+1})
    \bigr]_+.
\]
Оно сохраняет \(\lambda^k\ge0\). В оптимальной точке выполняются
\[
    g_j(x^\star)\le0,
    \qquad
    \lambda_j^\star\ge0,
    \qquad
    \lambda_j^\star g_j(x^\star)=0.
\]
Если ограничение неактивно, то \(g_j(x^\star)<0\) и обязательно
\(\lambda_j^\star=0\). При активном ограничении множитель может быть
положительным или нулевым; в устном пояснении лекции эти случаи были
перепутаны.

\section{Реализация и двойственная интерпретация}

Практический алгоритм имеет внешний и внутренний циклы. При
фиксированных \(\nu^k\), \(\lambda^k\) и \(c_k\) приближённо
минимизируется дополненная функция, причём предыдущая прямая точка
используется как тёплый старт. Для равенств контролируются прямая и
стационарная невязки
\[
    r_{\mathrm p}^k=h(x^k),
\]
\[
    r_{\mathrm s}^k
    =
    \nabla f(x^k)
    +
    J_h(x^k)^{\mathsf T}\nu^k.
\]
При неравенствах дополнительно проверяются допустимость,
неотрицательность множителей и дополнительная нежёсткость
\[
    \lVert\lambda^k\circ g(x^k)\rVert
    \le
    \varepsilon_{\mathrm c}.
\]

Если
\[
    g_c(\nu)
    =
    \inf_x\mathcal L_c(x,\nu),
\]
то при единственности внутреннего минимизатора
\[
    \nabla g_c(\nu)=h(x(\nu)).
\]
Поэтому обновление
\[
    \nu^{k+1}
    =
    \nu^k+ch(x^{k+1})
\]
можно рассматривать как градиентный подъём по дополненной двойственной
функции с шагом \(c\). Квадратичное дополнение сглаживает двойственную
задачу, но может связать блоки прямой переменной. Последовательное
разрешение этих блоков приводит к ADMM.
